\begin{table}[htbp]
    \caption{Limites do valor ótimo de $Z$ por conjunto --- instância grande}
    \label{tab:res-limites-grande}
    \centering
    \begin{tabular}{crrrr}
        \hline
        & \multicolumn{3}{c}{\textbf{Gurobi}} & \textbf{Limite combinatório} \\
        \cline{2-4} \cline{5-5}
        \textbf{Conjunto} & \textbf{Limite superior} & \textbf{Limite inferior} & \textbf{\textit{Gap}} & \textbf{$LB_{comb}$} \\
        \hline
        01 & 811\,250,0 & 0,0 & --- & 30\,200,0 \\
        02 & 819\,575,0 & 0,0 & --- & 40\,800,0 \\
        03 & 820\,500,0 & 47\,700,0 & 94,2\% & 47\,300,0 \\
        04 & 830\,875,0 & 0,0 & --- & 15\,500,0 \\
        05 & 816\,925,0 & 0,0 & --- & 32\,700,0 \\
        06 & 815\,200,0 & 29\,650,0 & 96,4\% & 28\,600,0 \\
        07 & 821\,300,0 & 0,0 & --- & 13\,900,0 \\
        08 & 835\,450,0 & 0,0 & --- & 42\,100,0 \\
        09 & 809\,600,0 & 0,0 & --- & 20\,000,0 \\
        10 & 822\,025,0 & 0,0 & --- & 30\,300,0 \\
        \hline
    \end{tabular}
    \par\smallskip
    {\footnotesize Nota - --- indica que o \textit{solver} não obteve limite inferior não trivial dentro do teto de tempo. \textit{Gap} na convenção do \textit{solver}. $LB_{comb}$: limite inferior combinatório, conforme a Equação~\eqref{eq:limite-combinatorio}.}
\end{table}

\begin{table}[htbp]
    \caption{Resultados da heurística por conjunto --- instância grande (10 execuções por conjunto; $Z$ sem a penalidade de capacidade)}
    \label{tab:res-z-grande}
    \centering
    \begin{adjustbox}{max width=\textwidth}
    \begin{tabular}{crrrrrcr}
        \hline
        \textbf{Conjunto} & \textbf{$Z$ da execução de menor $F$} & \textbf{$V_{cap}$} & \textbf{$Z$ médio} & \textbf{CV} & \textbf{Melhor $Z$ viável} & \textbf{Execuções viáveis} & \textbf{Excesso máx.} \\
        \hline
        01 & 50\,050,0 & 0 & 51\,685,0 & 1,5\% & 50\,050,0 & 10/10 & 65,7\% \\
        02 & 62\,500,0 & 0 & 63\,560,0 & 1,5\% & 62\,500,0 & 10/10 & 53,2\% \\
        03 & 67\,700,0 & 0 & 68\,612,5 & 0,9\% & 67\,700,0 & 10/10 & 41,9\% \\
        04 & 37\,825,0 & 0 & 39\,185,0 & 1,7\% & 37\,825,0 & 10/10 & 144,0\% \\
        05 & 52\,550,0 & 0 & 53\,830,0 & 1,3\% & 52\,550,0 & 10/10 & 60,7\% \\
        06 & 46\,700,0 & 0 & 48\,522,5 & 1,7\% & 46\,700,0 & 10/10 & 57,5\% \\
        07 & 35\,450,0 & 0 & 37\,377,5 & 2,5\% & 35\,450,0 & 10/10 & 155,0\% \\
        08 & 60\,975,0 & 0 & 62\,785,0 & 1,5\% & 60\,975,0 & 10/10 & 44,8\% \\
        09 & 35\,975,0 & 0 & 37\,415,0 & 2,0\% & 35\,975,0 & 10/10 & 79,9\% \\
        10 & 49\,650,0 & 0 & 51\,357,5 & 1,9\% & 49\,650,0 & 10/10 & 63,9\% \\
        \hline
    \end{tabular}
    \end{adjustbox}
    \par\smallskip
    {\footnotesize Nota - $Z$ da execução de menor $F$: valor de $Z$ da execução com menor \textit{fitness}, conforme a Equação~\eqref{eq:fitness}. $V_{cap}$: violação de capacidade dessa execução, em pacientes-dia, conforme a Equação~\eqref{eq:violacao-capacidade}, sendo $V_{cap} > 0$ uma solução inviável para o modelo exato. CV: coeficiente de variação de $Z$ entre as execuções (desvio padrão dividido pela média), como medida de dispersão. Melhor $Z$ viável: menor $Z$ entre as soluções viáveis encontradas nas execuções (--- se nenhuma). Excesso máx.: excesso máximo da execução de menor $F$ sobre o ótimo, conforme a Equação~\eqref{eq:excesso-otimo}, calculado com o maior entre o limite inferior do Gurobi e $LB_{comb}$.}
\end{table}
